\section{Velocidad angular y aceleración centrípeta}

\subsection{Ángulo, velocidad y aceleración angular}
\begin{align}
    \theta &= \frac{s}{r} \quad (\text{radianes}) \\
    \omega &= \frac{\Delta \theta}{\Delta t} \quad (\text{rad/s}) \\
    \alpha &= \frac{\Delta \omega}{\Delta t} \quad (\text{rad/s}^2)
\end{align}

\subsection{Relaciones lineales-angulares}
\begin{align}
    v &= \omega \cdot r \\
    a_t &= \alpha \cdot r \quad (\text{aceleración tangencial}) \\
    a_c &= \frac{v^2}{r} = \omega^2 r \quad (\text{aceleración centrípeta})
\end{align}

\subsection{Periodo y frecuencia}
\begin{align}
    T &= \frac{1}{f} \quad (\text{periodo}) \\
    \omega &= \frac{2\pi}{T} = 2\pi f
\end{align}


\resumebox{ejercicio}{Fuerza magnética del protón}{\footnotesize
    Una rueda de \(0.5 \, \text{m}\) de radio gira a \(300 \, \text{rpm}\). Calcula:
    \begin{enumerate}
        \item La velocidad angular en rad/s.
        \item La velocidad lineal en el borde de la rueda.
        \item La aceleración centrípeta.
    \end{enumerate}
}